\section{A mixed trace inequality}\label{sec:trace}

The triangular form is naturally represented by a cyclic product of
operators. Such a representation also appears in the dyadic model of
Kova\v{c}, Thiele and Zorin-Kranich~\cite[Equation~(1.5)]{KovacThieleZorinKranich2015}.
We prove a finite-dimensional estimate that controls this product by
three of the quadratic costs from Section~\ref{sec:matrix}.

Let $H_v$, $v\in\mathbb Z/3\mathbb Z$, be finite-dimensional complex
Hilbert spaces, and let $W_v:H_{v+1}\to H_v$ be arbitrary linear maps.
Combine the two edges incident to $v$ into the \emph{row star}
\begin{equation}\label{eq:row-star}
 R_v=[\,W_v\quad W_{v-1}^*\,]:
 H_{v+1}\oplus H_{v-1}\longrightarrow H_v.
\end{equation}
Its two Gram matrices are
\begin{equation}\label{eq:star-grams}
 T_v=R_v^*R_v,\qquad S_v=R_vR_v^*.
\end{equation}
Their positive eigenvalues are the squares of the nonzero singular
values of $R_v$, with the same multiplicities. In particular,
$\Phi(T_v)=\Phi(S_v)$.
Figure~\ref{fig:triangle} shows the cyclic edges and the row star at
vertex $0$.
\begin{figure}[htbp]
\centering
\begin{tikzpicture}[>=Stealth, every node/.style={font=\small},
 space/.style={circle,draw=black!60,minimum size=9mm,inner sep=2pt},
 edge/.style={->,thick,blue!65!black}]
 \node[space] (h0) at (0,1.9) {$H_0$};
 \node[space] (h1) at (-1.8,-.8) {$H_1$};
 \node[space] (h2) at (1.8,-.8) {$H_2$};
 \draw[edge] (h1) -- node[left=5pt] {$W_0$} (h0);
 \draw[edge] (h2) -- node[below=5pt] {$W_1$} (h1);
 \draw[edge] (h0) -- node[right=5pt] {$W_2$} (h2);
 \node[below=18pt] at (0,-.8) {the cyclic product $W_0W_1W_2$};
 \node[space] (r1) at (5,1.3) {$H_1$};
 \node[space] (r2) at (5,-.8) {$H_2$};
 \node[space] (r0) at (7.5,.25) {$H_0$};
 \draw[edge] (r1) -- node[above=5pt] {$W_0$} (r0);
 \draw[edge] (r2) -- node[below=5pt] {$W_2^*$} (r0);
 \node[below=18pt] at (6.25,-.8) {$R_0=[\,W_0\quad W_2^*\,]$};
\end{tikzpicture}
\caption{The three edge maps and the row star at vertex $0$.
Taking the adjoint of $W_2$ makes both blocks of $R_0$
have target $H_0$. The two Grams $R_0^*R_0$ and $R_0R_0^*$
have the same positive eigenvalues on different spaces.}
\label{fig:triangle}
\end{figure}


\begin{lemma}[Mixed trace estimate]\label{lem:mixed-trace}
Let $W_v:H_{v+1}\to H_v$ be linear maps between finite-dimensional
complex Hilbert spaces, with row stars $R_v$ and Gram matrices $T_v$
as in \eqref{eq:row-star}--\eqref{eq:star-grams}.
For Hermitian maps $D_v:H_v\to H_v$, put
$D_* =\max_v\lVert D_v\rVert_{\op}$. Then
\begin{equation}\label{eq:mixed-trace}
 \left|\tr(D_0W_0D_1W_1D_2W_2)\right|
 \le \frac53D_*\sum_{v=0}^2
 b_{T_v}(R_v^*D_vR_v).
\end{equation}
The constant is independent of the dimensions of the spaces.
\end{lemma}

\begin{proof}
We first put the three quadratic costs in a common metric. We then
estimate a cubic trace in an eigenbasis and recover the absolute value
of the original cyclic product by changing the phase of one edge.

\paragraph{A common quadratic form.}
On $H=H_0\oplus H_1\oplus H_2$, define the Hermitian matrices
\[
 M=\begin{pmatrix}
 0&W_0&W_2^*\\
 W_0^*&0&W_1\\
 W_2&W_1^*&0
 \end{pmatrix},
 \qquad D=\diag(D_0,D_1,D_2).
\]
Thus $\lVert D\rVert_{\op}=D_*$. Let $P_v$ be the orthogonal
projection onto $H_v$ and set
\[
 A_v=MP_vM,\qquad K_v=MP_vDM.
\]
Only the $v$-th block row of $P_vM$ can be nonzero; it consists of
$R_v$ with its columns placed in the corresponding vertex spaces.
Consequently $A_v$ and $K_v$ are,
respectively, $T_v$ and $R_v^*D_vR_v$, after a permutation of blocks
and extension by zero. These operations preserve the quadratic cost.
Since $P_v$ commutes with $D$, the matrix $K_v$ is Hermitian. Moreover,
\[
 \ker A_v=\ker(P_vM)\subseteq\ker K_v,
 \qquad 0\le A_v\le M^2.
\]
The first identity follows from
$\langle A_vx,x\rangle=\lVert P_vMx\rVert^2$.
Thus $K_v$ is supported on $A_v$, and hence also on $M^2$.
Lemma~\ref{lem:order}, followed by the quadratic-form inequality for
three vectors, yields
\begin{equation}\label{eq:trace-common-metric}
 \begin{aligned}
 \mathcal E
 &:=\sum_v b_{T_v}(R_v^*D_vR_v)\\
 &\ge\sum_v b_{M^2}(K_v)
 \ge\frac13 b_{M^2}\left(\sum_vK_v\right)
 =\frac13 b_{M^2}(MDM).
 \end{aligned}
\end{equation}
All these quadratic forms are taken on their supports; no
invertibility of $M$ is required.

\paragraph{The cubic trace.}
Choose an orthonormal eigenbasis $(e_i)$ of $M$, write the real
eigenvalues as $m_i$, and arrange the indices so that $q_i:=|m_i|$ is
nonincreasing. Equal values may be ordered arbitrarily. Write $D_{ij}$
for the entries of $D$ in this basis. The spectral formula for
$b_{M^2}$ gives
\begin{equation}\label{eq:trace-Q}
 \mathcal E\ge\frac12Q,
 \qquad
 Q:=\sum_{i,j:q_i+q_j>0}
 \frac{q_i^2q_j^2}{q_i+q_j}|D_{ij}|^2.
\end{equation}
If either eigenvalue is zero, its summand is zero; pairs of zero
eigenvalues are omitted. Define
\[
 A_i=q_i^3|D_{ii}|^2,\qquad
 E_i=q_i\sum_{j>i}q_j^2|D_{ij}|^2,
 \qquad A=\sum_iA_i,\quad E=\sum_iE_i.
\]
Hermitian symmetry and the ordering of the $q_i$ imply
\begin{equation}\label{eq:trace-AE}
 \begin{aligned}
 Q
 &=\frac A2+
 2\sum_{\substack{i<j\\q_i+q_j>0}}
 \frac{q_i^2q_j^2}{q_i+q_j}|D_{ij}|^2
 \ge\frac A2+E.
 \end{aligned}
\end{equation}
Indeed, for $i<j$ we have $q_j\le q_i$, so
$2q_i/(q_i+q_j)\ge1$ whenever the denominator is positive.

We now bound
\[
 \tr((DM)^3)
 =\sum_{i,j,k}m_im_jm_kD_{ij}D_{jk}D_{ki}
\]
by grouping the triples according to how often their smallest index
occurs. This is the index of a largest absolute eigenvalue in the
triple.

If the smallest index $i$ occurs once, its three cyclic positions
give
\[
 3m_i\sum_{j,k>i}m_jm_kD_{ij}D_{jk}D_{ki}.
\]
Here $j=k$ is allowed. For the vector $x=(m_jD_{ji})_{j>i}$, the inner
sum is $x^*D_{>i}x$, where $D_{>i}$ is the compression of $D$ to the
indices greater than $i$. Since
$\lVert D_{>i}\rVert_{\op}\le D_*$, the modulus of this contribution
is at most $3D_*E_i$. This argument retains the signs of the $m_j$
and uses $D_{ji}=\overline{D_{ij}}$.

If the smallest index $i$ occurs twice, the contribution is
\[
 3m_i^2D_{ii}\sum_{k>i}m_k|D_{ik}|^2.
\]
Cauchy--Schwarz and
$\sum_k|D_{ik}|^2=\lVert De_i\rVert^2\le D_*^2$ give
\[
 \left|\sum_{k>i}m_k|D_{ik}|^2\right|
 \le D_*\left(\sum_{k>i}q_k^2|D_{ik}|^2\right)^{1/2}.
\]
Its contribution therefore has modulus at most
$3D_*\sqrt{A_iE_i}$. This also holds when $q_i=0$, without any
division by $q_i$.

Finally, when all three indices equal $i$, the modulus is at most
$D_*A_i$. These three cases exhaust all triples, including repeated
indices. Summing them and applying Cauchy--Schwarz once more, we obtain
\begin{equation}\label{eq:trace-cubic}
 \begin{aligned}
 |\tr((DM)^3)|
 &\le D_*\bigl(A+3E+3\sqrt{AE}\bigr)\\
 &\le D_*\left(\frac52A+\frac92E\right)
 \le 5D_*\left(\frac A2+E\right)
 \le5D_*Q.
 \end{aligned}
\end{equation}

\paragraph{Recovering the complex cyclic product.}
We have bounded the cubic trace by the star costs. It remains to
extract the modulus of the desired triangular trace. Put
\[
 \tau=\tr(D_0W_0D_1W_1D_2W_2).
\]
Because $M$ has zero diagonal blocks, a nonzero closed three-step
block product must visit all three vertices. The three starting
points in one orientation each give $\tau$, by cyclicity of the
trace. Reversing the orientation gives $\overline\tau$, since the
$D_v$ are Hermitian. Hence
\[
 \tr((DM)^3)=6\operatorname{Re}\tau.
\]
If $\tau\ne0$, choose $z\in\mathbb C$ with $|z|=1$ and
$z\tau=|\tau|$, and replace $W_0$ by $zW_0$. The row stars change by
\[
 \begin{aligned}
 R_0&\longmapsto
 R_0\diag(zI_{H_1},I_{H_2}),\\
 R_1&\longmapsto
 R_1\diag(I_{H_2},\overline zI_{H_0}),
 \end{aligned}
\]
while $R_2$ is unchanged. Each displayed right multiplier is unitary
and conjugates the corresponding Gram and inserted Gram by the same
unitary. Thus $\mathcal E$ is unchanged, as is $D_*$.
Applying \eqref{eq:trace-Q} and \eqref{eq:trace-cubic} to the modified
matrices, with $Q$ computed for the new $M$, gives
\[
 6|\tau|\le5D_*Q\le10D_*\mathcal E,
\]
which is \eqref{eq:mixed-trace}. The case $\tau=0$ is immediate.
\end{proof}
